\section{Experiments}

\subsection{Questions}

Our evaluation is organized around five questions:
\begin{enumerate}
    \item Does an easy predictable tag dominate planning geometry despite
    healthy anti-collapse statistics?
    \item Does action supervision merely make physical information decodable,
    or does it also make that information load-bearing for planning?
    \item Are gains explained by negative-conflict removal, generic prediction
    attenuation, or sample-specific control direction?
    \item Can soft residual admission preserve Reacher dynamics without
    restoring watermark dominance on PushT?
    \item What happens when a static factor is decision-causal, as in RandGoal?
\end{enumerate}

\subsection{Tasks and nuisance construction}

We use PushT, RandGoal, TwoRooms, OGBench-Cube, and Reacher under the LeWM
training and latent-space planning framework.  The tagged condition adds a
small episode-constant visual watermark that is predictable but does not alter
the environment transition or expert decision.  Clean and tagged variants
share trajectories, optimization, architecture, and planner settings whenever
possible.

RandGoal is a deliberate counterexample to static-feature filtering.  Holding
physical state fixed while changing the goal changes the expert's future
actions; the goal is therefore decision-causal even when it remains constant
within an episode.  A random watermark changes neither action selection nor
transition dynamics.  A successful allocation method should suppress the
latter without erasing the former.

\subsection{Metrics and protocol}

Planning is evaluated with the same CEM configuration, candidate budget,
evaluation seeds, and checkpoint rule within each matched comparison.  We
report success rate over complete learning curves, a predeclared late-window
mean, and final performance; single 50-episode peaks are diagnostic only.
Selected comparisons use at least three training seeds and 200--500 evaluation
episodes per selected checkpoint.

Representation evaluation includes same-content and same-tag cosine
similarity, frozen probes for task variables and tag identity, prediction and
control losses, gradient cosine, negative-conflict fraction, and admitted
prediction-gradient norm.  The content--tag margin is
\begin{equation}
    m=\mathrm{cos}_{\mathrm{same\ content}}-
    \mathrm{cos}_{\mathrm{same\ tag}}.
\end{equation}
Positive $m$ indicates that task content, rather than watermark identity,
organizes the representation.  Probe accuracy is interpreted only as
decodability; it does not establish causal use by the planner.

\subsection{Matched methods}

We compare LeWM, identical-head IDM/control-only training, summed prediction
and control losses, zero-floor cosine \shortmethod{}, soft admission with
$\beta\in\{0.25,0.5,0.75\}$, the conflict-only comparator in
Eq.~\ref{eq:pcgrad-comparator}, norm-matched scalar attenuation, and a
retention-matched shuffled guide.  State alignment is discussed as a
privileged-supervision reference but is not included as a like-for-like
training baseline.

\subsection{The orthogonality challenge}
\label{sec:orthogonality-challenge}

The central methodological objection is that near-orthogonality is common in
high dimension.  A small value of $g_p^\top g_c$ therefore does not by itself
show that the corresponding predictive update is task-irrelevant.  It can
instead mean that the control guide is low-rank, that a useful control feature
is already decoded with small residual error, that different control
components cancel in their sum, or that prediction supplies dynamics needed
only over a longer planning horizon.  We consequently do not treat
orthogonality as a semantic label.  We test what operation on the orthogonal
component is actually responsible for the planning result.

Our identification strategy uses five matched interventions.  First,
\emph{parallel-only hard drop} removes the entire orthogonal component; if it
underperforms cosine admission, orthogonal predictive information is useful.
Second, the conflict-only comparator in
Eq.~\ref{eq:pcgrad-comparator} removes only genuinely negative interference; if
it behaves like LeWM while cosine admission improves tagged planning, rare
negative conflicts are insufficient to explain the gain.  Third,
norm-matched scalar attenuation tests whether the result is merely a lower
encoder-side prediction learning rate.  Fourth, soft admission sweeps the
orthogonal residual floor and traces the trade-off between nuisance resistance
and preservation of unknown dynamics.  Fifth, the component-span diagnostic
tests whether the summed control gradient appears orthogonal because distinct
horizon, inverse, cycle, and reachability gradients cancel.

These interventions yield a falsifiable interpretation.  Evidence for
directional control guidance requires cosine admission to outperform both
norm-matched attenuation and a retention-matched shuffled guide.  Evidence
that orthogonal information is sometimes necessary requires hard drop to hurt
a clean dynamics task and soft admission to repair it.  Evidence that the
watermark effect is not standard gradient-conflict resolution requires the
conflict-only rule to fail to reproduce the tagged-task gain.  We report raw
cosine and admitted-norm distributions alongside success rates so that a
method cannot appear successful merely by silently suppressing most encoder
learning.

\subsection{Predeclared interpretation}

Soft admission is considered successful only if one shared setting improves
Reacher while retaining the PushT watermark benefit; choosing a different
$\beta$ per task is not evidence for a general method.  Conflict-only routing
tests whether negative interference suffices.  The component-span diagnostic
tests the high-dimensional orthogonality objection.  Multi-seed experiments
are promoted only after these seed-zero mechanism checks pass.
